% problem-000253 4 N-H键活化与金属氢化物
\section*{4. N-H键活化与金属氢化物}
\textit{下列为分问。}

\subsection*{4-1}
过渡金属配合物对氨分子中N-H键的活化，通常包括过渡金属中心的氧化加成及其氢化物的去质子化两个步骤。科学家发展了一种金属铱配合物A活化N-H键。其结构如下：

C
B
（图：）
A

（图：）

（图：）

（图：）

\subsection*{4-1}
-1 分子A 的结构为平面四边形，可与 $\mathrm { N H } _ { 3 }$ 在室温下反应实现对N-H键的活化，金属氧化数发生了变化。确定反应前后金属Ir的氧化数，并画出产物的结构式。

\subsection*{4-1}
-2 B 也可以活化 $\mathrm { N H } _ { 3 }$ 中N-H键。确定产物中 Zr的氧化数，判断反应前后金属 Zr的氧化数是否有变化。

\subsection*{4-2}
除了过渡金属配合物外，主族化合物C也能活化 $\mathrm { N H } _ { 3 }$ 中N-H键。

\subsection*{4-2}
-1画出当P的形式电荷为-1时化合物C所有稳定的共振结构，并标出所有形式电荷。

\subsection*{4-2}
-2画出当P的形式电荷为-2时化合物C所有稳定的共振结构，并标出所有形式电荷。

\subsection*{4-2}
-3C与NH3反应后的产物为D。D中P的氧化数为+5，写出D中P原子的配位几何结构。

\subsection*{4-3}
过渡金属配合物活化 N-H 键时，一般不会生成 $\mathrm { H } _ { 2 }$ ，这是因为气态 $\mathrm { N H } _ { 3 }$ 分子中N-H键的解离能高达416kJ $\mathrm { \ m o l ^ { - 1 } }$ 。当 $\mathrm { N H } _ { 3 }$ 分子和过渡金属配位时，N-H键虽然被削弱，但常见的Co(II)、Ru(II)和Pt(II)等过渡金属氨配合物对N-H键的削弱程度有限，使其从热力学上无法自发形成 $\mathrm { H } _ { 2 \mathrm { ~ c ~ } }$ 然而，有些不常见的金属氨配合物能在极大程度上削弱N-H键，如含配体 $\mathsf { P h } _ { \mathrm { T p y } }$ 的金属M氨配合物，其合成过程如下：

\subsection*{4-3}
-1 将 $\mathsf { P h } _ { \mathrm { T p y } }$ 与 $\mathbf { M } ( \mathrm { T H F } ) _ { 3 } \mathbf { C } \mathrm { l } _ { 3 }$ 在150 mL THF 溶液中反应 24 小时，生成产物E。经纯化干燥后，E的元素分析结果(质量百分比)为： $\mathbf { C } ,$ 49.29; H, 2.95; N, 8.21。通过计算判断金属M为何种元素。

\subsection*{4-3}
-2 在Na/Hg齐存在下，将0.395g(1.974mmol)配体 $\mathbf { M e P P h } _ { 2 }$ 与含0.500g(0.977mmol) E的 THF 溶液反应，生成产物 G。G 经纯化结晶得 0.429 g，有效磁矩为$1 . 7 \mu 8$ 。写出产物G的分子式，确定金属中心的价电子构型，并计算反应产率。4-3-3G经如下反应转化为含氨作为配体的化合物H，升高反应体系温度形成化合物I，并释放出 $\mathrm { H } _ { 2 } ;$

$
\mathrm{G} \xrightarrow [ \mathrm{PhH} , \mathrm{rt} , - \mathrm{NaCl} ]{\mathrm{Na} (\mathrm{B} (3 , 5 - \mathrm{CF} _ {3}) \mathrm{C} _ {6} \mathrm{H} _ {3}) _ {4} , 1 \text {equiv.} \mathrm{NH} _ {3}} \mathrm{H} \xrightarrow {\mathrm{PhH} , 6 0 ^ {\circ} \mathrm{C}} \mathrm{I} + 1 / 2 \mathrm{H} _ {2}
$

写出配合物H和Ⅰ阳离子的化学式，确定中金属中心的氧化数。

\subsection*{4-3}
-4为了解配体对配合物H中N-H键强度的影响，科学家进行了如下的实验(分子H以 $[ \mathbf { M } \mathbf { - N H } _ { 3 } ]$ +表示):

在 THF 溶剂中，测得电极反应 $\mathrm { \small { [ M - N H _ { 3 } ] ^ { ( m + 1 ) + } + e ^ { - } = [ M - N H _ { 3 } ] ^ { m + } } }$ 的标准电极电势 $\varphi _ { 1 } ^ { \ominus } = - 0 . 5 5 6 \mathrm { V }$ (相对于THF 溶剂中标准氢电极)；在 THF 溶剂中，测定[M$\mathrm { N H } _ { 3 } ] ^ { ( \mathrm { m + 1 } ) + }$ 的 $\mathtt { p } K _ { \mathtt { a } } = 3 . 6$ 。气态氢原子的标准生成吉布斯自由能为203.3kJmol-1，气态氢原子在THF中的溶剂化摩尔吉布斯自由能为21.42kJ $\mathrm { \ m o l ^ { - 1 } }$ 。计算以THF为溶剂时，H中N-H键的解离吉布斯自由能。
